% Coupling speciation to competition --- plan \S3.3, "\S4".
% Sources: 09/05:206-323, minus the block quotation at :213-290 (parked as
% _parked/nietzsche_bge_262.tex) and minus the confession at :321-323 (cut;
% replaced by the "What this chapter is" paragraph of \S1).
% Renames beyond plan \S7.2, both forced by \S7.1's rule against contradicting
% a thesis-wide symbol: the fitness increment \delta -> \varepsilon (\delta is
% the catastrophe rate, and this quantity is the evolution rate of the
% preceding section under a second letter), and the speciation exponent
% \lambda -> \zeta (\lambda is the birth rate).  See PHASE_A_REPORT.md 9.1.
% The source at 09/05:276,278 prints both step probabilities with the same
% left-hand side, \Pr(\Delta A_t = 1); the second is written here as -1.

\subsection{Coupling speciation to competition}
\label{var:sec:pimentelplus}

\Cref{var:sec:pimentel} makes a losing population strong, the squeeze being
what drives its resistance up. This subsection asks the complement, whether a
winning population can be made weak by a mechanism of the same kind; the
hypothesis it states is that excess variation weakens an established type, so
that a population which is winning diversifies and thereby loses. The claim
is Nietzsche's, put into a model: a type is made hard by long struggle
against constant unfavourable conditions, favourable conditions relax the
constraint which made it hard, and the resulting exuberance of variation
is itself the mechanism of decline \cite{nietzsche2013beyond}. What is worth retaining is the shape of the claim: that the decline is
\emph{endogenous}, a consequence of the success preceding it rather than a
response to any external shock.

The model is the kernel of \cref{var:sec:pimentelmodel}, recast so that the
quantity under selection is fitness rather than resistance, with a species
count carried alongside and allowed to act back on that fitness. Written as
step probabilities for $A_t$,
\begin{equation}
\pr{\Delta A_t=1}=\frac{A_t}{N}\,\frac{f_a}{\bar f},
\qquad
\pr{\Delta A_t=-1}=\lrb{1-\frac{A_t}{N}}\frac{f_b}{\bar f},
\label{var:eq:fitkernel}
\end{equation}
with $\bar f$ the mean fitness across the domain and, under the Moran
assumption, $A_t=N-B_t$. Each probability is again the chance of meeting a
member of the opposing group multiplied by the chance of prevailing against
them, as in \cref{var:eq:kernel}; what has changed is that prevailing is now
governed by a fitness ratio rather than by a blocking probability. Fitness is
driven up by minority status exactly as resistance was:
\begin{equation}
f_a(t+\Delta t)=f_a(t)+\varepsilon\lrb{\frac{B_t}{A_t}}^{\!\alpha},
\qquad
f_b(t+\Delta t)=f_b(t)+\varepsilon\lrb{\frac{A_t}{B_t}}^{\!\alpha}.
\label{var:eq:fitupdate}
\end{equation}

Fitness here is the same evolving quantity as the resistance $R_a,R_b$ of
the previous section, under a second name: the update rule and its exponent
are identical, and the increment coefficient is $\varepsilon$ in both. The
two kernels are not literally identical, since \cref{var:eq:kernel} has a
stalemate outcome and \cref{var:eq:fitkernel} has none; both are carried
because the fitness form is the one that admits the coupling below, in which
a scalar can be added to and subtracted from.

Now give each population its own species count, $S^{(a)}_t$ and $S^{(b)}_t$,
evolving as birth--death processes with constant per-species extinction
$\mu_a$ and a speciation rate that rises with both dominance and fitness:
\begin{equation}
\pr{\Delta S^{(a)}_t=1}=\beta_a S^{(a)}_t\,\Delta t,
\qquad
\pr{\Delta S^{(a)}_t=-1}=\mu_a S^{(a)}_t\,\Delta t,
\qquad
\beta_a=\lrb{\frac{A_t}{N}}^{\!\zeta}f_a,
\label{var:eq:subspec}
\end{equation}
and correspondingly for $b$ with $A_t$ replaced by $B_t$, the exponent
$\zeta$ controlling how firmly dominance must be established before variation
is released. This is the premise of \cref{var:sec:speciation} transplanted
into a competitive setting: variation is in constant supply, and what
releases it is the easing of whatever was suppressing it, measured here by
the share of the domain a population controls.

The coupling that closes the loop is a penalty on fitness proportional to the
species count:
\begin{equation}
f_a(t+\Delta t)=f_a(t)+\varepsilon\lrb{\frac{B_t}{A_t}}^{\!\alpha}
-\gamma S^{(a)}_t,
\qquad
f_b(t+\Delta t)=f_b(t)+\varepsilon\lrb{\frac{A_t}{B_t}}^{\!\alpha}
-\gamma S^{(b)}_t,
\label{var:eq:fitpenalty}
\end{equation}
with $\gamma$ the amount by which the inefficiency of carrying excess
variation weighs on competitive performance.

That is the whole specification, and it has never been simulated. Both routes
to reversal are present in it at once: the squeeze that strengthens the
loser, through \cref{var:eq:fitupdate}, and the variation that weakens the
winner, through the penalty of \cref{var:eq:fitpenalty}. Both favour whichever
population is behind, so the model does not say in advance which of them does
the work, and that is what makes it worth running.

\paragraph{Open questions.}\label{var:rem:open-plus}
\begin{enumerate}[leftmargin=1.6em,itemsep=0.2em,topsep=0.2em]
\item \textbf{Whether collapse by variation happens at all.} Does the
$-\gamma S^{(a)}_t$ term in \cref{var:eq:fitpenalty} ever produce a reversal
that the squeeze-driven mechanism of \cref{var:sec:pimentel} would not have
produced unaided, or does it only slow that one down? Separating the two
means holding everything else fixed and comparing runs at $\gamma=0$ against
runs at $\gamma>0$, which asks for simulation and for nothing else.
\item \textbf{The relevant regime.} Collapse by variation should require
$\gamma S^{(a)}_t$ to grow comparable with the accumulated fitness advantage
over the timescale on which $S^{(a)}_t$ itself grows, which ties $\gamma$,
$\zeta$ and $\mu_a$ together. Whether that regime is a thin sliver of
parameter space or a broad band of it follows directly from the first
question.
\end{enumerate}
